\section{Related work}
\label{sec:related}

\paragraph{Ray tracing and path tracing.} Whitted's recursive ray tracing
\cite{whitted1980improved} introduced shadows, reflection, and refraction by
spawning secondary rays; Kajiya's rendering equation \cite{kajiya1986rendering}
cast image synthesis as an integral over all light paths, solved by Monte-Carlo
path tracing. Our tracer is a minimal path tracer (direct sun plus one diffuse
bounce), enough to exhibit both the geometric cost that acceleration structures
attack and the stochastic noise that sample budgets and denoisers attack. The
ray/triangle test is Möller--Trumbore \cite{moller1997fast}; the standard
reference for the whole pipeline is \emph{Physically Based Rendering}
\cite{pharr2016physically}.

\paragraph{Acceleration structures.} The idea that a ray should not be tested
against every primitive dates to the earliest complex-scene renderers
\cite{rubin1980representation}. Bounding-volume hierarchies are now the dominant
structure, including for dynamic scenes \cite{wald2007ray} and in hardware
ray-tracing cores. We use a textbook median-split BVH; our interest is not in a
better build but in \emph{measuring} what it buys.

\paragraph{Instancing.} Storing one mesh and referencing it through many
transforms --- a two-level hierarchy with a bottom-level structure (BLAS) per
mesh and a top-level structure (TLAS) over instances --- is standard in production
and is the mechanism behind dense vegetation such as the grass of \emph{Breath of
the Wild}. It is lossless by construction; we verify this and quantify its cost
and memory savings.

\paragraph{Level of detail.} Replacing distant geometry with cheaper
approximations is a classic real-time technique \cite{luebke2003level}; billboard
impostors are its aggressive limit. Unlike the lossless methods, LOD changes the
image, and the quality/cost trade is exactly what we measure.

\paragraph{Sampling and denoising.} Monte-Carlo estimators converge as
$N^{-1/2}$ in the number of samples, so halving noise costs four times the rays.
Learned and hand-built denoisers \cite{chaitanya2017interactive} make low-sample
renders usable and underpin production systems (e.g.\ NVIDIA OptiX, DLSS). Our
denoiser is a cross-bilateral filter \cite{tomasi1998bilateral} guided by
G-buffer normals and depth --- the minimal version of that idea, chosen so its
behaviour is transparent rather than opaque.
